\scanpage{73}
dans $(\mathrm{Ens})$ qui commutent aux $\varprojlim$ sont les foncteurs représentables) i.e. si $T=T_R$, on a $R=S^\circ$, où $S=T(\mathrm{Ens})$. \deleted{Cela précise} \deleted{Donc, \unsure{pour les}} Ou encore $R^\circ=S$. Donc :

\uline{Cor} \deleted{les} \inserted{2-cat des} structures algébriques définissables par $\varprojlim$ (sans axiome d'exactitude) \deleted{correspondent aux catégories} est 2-équiv. par $T\longmapsto T(\mathrm{Ens})$ à celle des $\mathcal U$-catégories $S$ avec $\varinjlim$ quelc. et \inserted{petites} \deleted{familles} \inserted{sous-cat} génératrices par \unsure{épis.|épi.} stricts [i.e. telles que $S\longrightarrow\widehat{I^{\unsure{a|\circ}}}$ soit conservatif]. La donnée d'une base \inserted{de $T$} équivaut à celle d'une famille d'objets cogénératrice de $S$. Pour toute $C\in\mathrm{Ob}\,\Delta$, on a
\[
T(C)\simeq\underline{\mathrm{Hom}}_{\Delta}(S^\circ,\unsure{\mathscr C|\mathbf C|C})\deleted{\unsure{\mathrm{Ens}}}
\]

L'équivalence est obtenue en notant que le deuxième mb est formé des foncteurs tels que $\forall x\in C$, le composé avec $\underline{\mathrm{Hom}}_{C}(x,-)$ est représentable, donc le deuxième mb \deleted{est} peut s'interpréter comme les foncteurs $S^\circ\times C^\circ\longrightarrow(\mathrm{Ens})$ qui sont représentables pour chaque argument fixé, ou encore les foncteurs $C^\circ\longrightarrow S$ qui sont tels que pour tout $Y\in S$, son composé avec $\mathrm{Hom}(Y,-)$ soit représentable.

Or on a
\[
T(C)\longrightarrow\underline{\mathrm{Hom}}(C^\circ,\underbrace{T(\mathrm{Ens})}_{\substack{\Vert\\S}})=\underline{\mathrm{Hom}}(\unsure{C^\circ|\mathscr C^\circ},S)
\]
déduit de $C^\circ\longrightarrow\underline{\mathrm{Hom}}_{\Delta}(C,\mathrm{Ens})\subset\widehat{C^\circ}$ induisant $C^\circ\longrightarrow\underline{\mathrm{Hom}}(T(C),T(\mathrm{Ens}))$ i.e. $T(C)\longrightarrow\deleted{\unsure{\mathrm{Hom}|T(C)}}\underline{\mathrm{Hom}}(C^\circ\ T(\mathrm{Ens}))$.

\inserted{\deleted{\unsure{Ens}} \textcircled{4} Comme $S=T(\mathrm{Ens})$ a des $\varprojlim$ quelc, et $S\cong R^\circ$, $R$ a aussi des $\varinjlim$ quelc, donc $S$ et $R$ ont toutes les $\varinjlim$ et $\varprojlim$}

\textcircled{\deleted{4}5} Soit un foncteur $\varphi:R\longrightarrow R'$ entre objets de $\Delta$. Si $R$ a une petite famille cogénératrice, alors $\varphi$ est un $\Delta$-foncteur, i.e. commute aux $\varprojlim$, sss il admet un adjoint à gauche $\psi$
\[
\mathrm{Hom}(y',\varphi(y))\simeq\mathrm{Hom}(\psi(y'),x)
\]
% slanted
\margin{ainsi les \inserted{$\Delta$-}foncteurs $\varphi:R\longrightarrow R'$ (i.e. $T'\longrightarrow T$) correspondent aux foncteurs $\psi:R'\longrightarrow R$ (i.e. $\psi^\circ:S'\longrightarrow S$) qui ont un adjoint à droite (\unsure{Corresp à|corresp. à} \unsure{$\varphi^\circ$})}
Alors $\varphi$ est un passage à une sous-cat. \unsure{mise|univers.} en fractions sss \deleted{$R'$ est une sous-catégorie} $\psi$ est pl fidèle, \inserted{(et alors $R'$ a une p. famille cog.)} Ainsi les \uline{sous-}\deleted{catégories}\inserted{th$^{\text{é}}$}théories \inserted{Alg.} de la th. algébrique $R$ correspondent \deleted{aux sous-} à \uline{certaines} sous-catégories pleines de $R$, savoir celles qui ont des $\varprojlim$ et telles que les foncteurs d'inclusion $\psi$ ait un adjoint à \deleted{\unsure{g}} droite. En termes de $S=R^\circ=T(\mathrm{Ens})$, et \inserted{$S'=R'^\circ=T'(\mathrm{Ens})$} \deleted{ce sont les sous-catégo-}
% archive sheet number: 72
